%% =================================================================
%% Beber, Lamon, Saveriano, Fontanelli, Palopoli.
%% "Autonomous Robotic Tissue Palpation and Abnormalities
%%  Characterisation via Ergodic Exploration."
%% IEEE Robotics and Automation Letters, preprint accepted Feb. 2026.
%% DOI: 10.1109/LRA.2026.3673907
%%
%% Reconstruction of page 6 as study notes.
%% Contents: Figure 3 (2x5 grid + 2x1 + 2x1 comparison panels),
%% continuation of section IV-B, section IV-C (BO comparison),
%% start of section IV-D (Segmentation, cuts off).
%% No numbered equations on this page.
%% =================================================================

\documentclass[11pt]{article}

\usepackage[margin=1in]{geometry}
\usepackage{amsmath, amssymb}
\usepackage{mathrsfs}
\usepackage{bm}
\usepackage{xcolor}
\usepackage{fancyhdr}
\usepackage{url}

\pagestyle{fancy}
\fancyhf{}
\lhead{\small L.~Beber et al.}
\rhead{\small IEEE RA-L Preprint, Accepted Feb.\ 2026}
\cfoot{\small 6 $\to$ \arabic{page}}
\renewcommand{\headrulewidth}{0.4pt}

\setcounter{page}{1}
\setcounter{equation}{6}   % page 5 had no numbered equations; counter still 6
\setlength{\parskip}{4pt}

\begin{document}

% ---------- Preprint notice strip ----------
\begin{center}
{\footnotesize\itshape
This article has been accepted for publication in IEEE Robotics and
Automation Letters. This is the author's version which has not been
fully edited and content may change prior to final publication.
Citation information: DOI 10.1109/LRA.2026.3673907\par}
\end{center}

\vspace{0.4em}
\hrule
\vspace{0.6em}

\noindent
{\footnotesize\scshape 6 \hfill IEEE Robotics and Automation Letters.\ Preprint Version.\ Accepted February, 2026}

\vspace{1.2em}

% ---------- Figure 3 placeholder ----------
\begin{center}
\fbox{\parbox{0.95\textwidth}{\centering\vspace{1em}
\textbf{[Figure 3 --- placeholder]}\\[0.8em]
A wide figure spanning the full page width, composed of three
sub-figures side by side. \textbf{Sub-figure (a)}
\emph{``Ergodic exploration with HEDAC (ours)''}: a 2-row $\times$
5-column mini-grid. The top row is labelled \textit{EID} (the
target Expected Information Density map) and the bottom row is
labelled \textit{Estimated Map}. The five columns are time snapshots
at $t = 14$, $28$, $42$, $56$, $70$~s, showing the progressive
evolution of both fields as the ergodic controller visits the
domain. \textbf{Sub-figure (b)} \emph{``BO-EI''}: a single 2-row
column showing the final EID and estimated map for the standard BO
baseline. \textbf{Sub-figure (c)} \emph{``BO-EIS''}: same layout for
the continuous-sampling BO variant. In all panels: a cyan circle
marks the start of the trajectory, a magenta square marks the end
of the trajectory, a red trace shows the collected sample points
along the executed path, and (for (b), (c) only) orange markers
indicate the discrete locations chosen by the BO acquisition
function. The overall visual impression in sub-figure (a) is that
the red trajectory gradually concentrates around the three hot-spot
centres as $t$ increases, matching the EID evolution.
\vspace{1em}}}
\end{center}

\vspace{0.3em}

\begin{center}
\textbf{Fig. 3.} End-effector tip trajectory during the exploration
of an elasticity distribution with three stiff regions: comparison
of different planners. The cyan circle indicates the beginning of
the trajectory, while the magenta square denotes the end of the
trajectory. In plots (b) and (c), the orange markers indicate the
locations selected by BO, whereas the red markers represent the
sampled points collected along the executed trajectory.
\end{center}

\vspace{1em}

% ---------- Continuation of IV-B Robustness from page 5 ----------
\noindent
time across the three scenarios, regardless of the presence of
measurement uncertainty. These results demonstrate the robustness of
the method with respect to both the number of stiff regions to be
detected and the influence of measurement noise. The RMSE, instead,
grows under both experimental conditions as the number of stiff
regions increases. This behaviour may be attributed to uncertainty
at the boundaries of the stiff regions, which are typically
under-sampled, and will be examined in greater detail in Sec.~IV-D.
As expected, we observe an increased RMSE in the second experimental
condition, attributable to the presence of the noise.

Figure~3(a) shows an example of the evolution of the agent
trajectory over time, when exploring the map with 3 stiff regions.
All stiffer regions are successfully detected within 40~s. Over
time, the ergodic metric $\alpha$ decreases from 0.79 at $t = 14$~s
to 0.38 at $t = 70$~s, reflecting improved coverage of the target
distribution. In parallel, the RMSE shows a substantial reduction,
dropping from 13.75~kPa to 2.85~kPa. Intermediate values reinforce
this trend: at $t = 28$~s, $\alpha = 0.54$ and
$\mathrm{RMSE} = 11.19$~kPa; at $t = 42$~s, $\alpha = 0.51$ and
$\mathrm{RMSE} = 3.7$~kPa; and at $t = 56$~s, $\alpha = 0.43$ and
$\mathrm{RMSE} = 2.56$~kPa. This consistent evolution confirms that
as the exploration becomes more ergodic, the accuracy of the
reconstructed stiffness map improves accordingly.

\vspace{0.6em}

% ---------- C. Comparison with BO Baselines ----------
\noindent\textit{C.~Comparison with BO Baselines}

\vspace{0.2em}

\noindent
In this section, we focus on the comparison of the results of our
ergodic approach with the two variants of BO described in
Sec.~IV-A. We consider the distribution with 3 stiff regions, which
has proven to be the most challenging for the ergodic exploration
(see Fig.~3). The simulation time is the sum of the simulated agent
motion with a velocity of 1~cm/s and the computation time required
to determine the next sampling location. Each simulation is
conducted under two conditions: one assuming noise-free elasticity
measurements, and another incorporating additive white noise
$\mathcal{N}(0,\,6.25\,\mathrm{kPa}^{2})$, similar to Sec.~IV-B.
Fixed lenght of 400, 600, and 800~mm, were used to evaluate the
performance evolution over time, ensuring a fair comparison
independently of the method-specific stopping criteria. The average
results of 100 simulations are listed in Tab.~II.

This comparison highlights that BO-EI consistently underperforms
compared to the other two methods, yielding the lowest DR and the
highest RMSE across all tested conditions. Both BO-EIS and the
ergodic strategy achieve comparable detection rates once sufficient
exploration has been performed, indicating successful identification
of the number and location of stiff regions. However, the ergodic
approach consistently provides lower reconstruction errors,
resulting in more accurate stiffness maps. At shorter travelled
distances (500~mm--650~mm), BO-EIS achieves higher detection rates,
reflecting faster initial identification of stiff regions. This
behaviour is expected given the exploitative nature of BO-based
strategies, which prioritise sampling high-uncertainty locations.
However, in clinical settings, once suspicious regions are
identified, accurate post-detection refinement becomes critical.
Therefore, it is more relevant to compare the approaches once
detection performance is maximised (800~mm, DR $\approx$ 99--100\%),
and differences between the methods are primarily reflected in
reconstruction accuracy and execution duration. In this regime, the
ergodic strategy achieves lower RMSE while requiring substantially
less execution time than BO-EIS under the same travelled-distance
budget, both in noise-free and noisy conditions. These results
indicate that, in the post-detection phase, ergodic exploration
enables more efficient refinement of stiffness maps, yielding
improved accuracy per unit motion and reduced execution time under
equal motion cost. Compared to BO-based approaches, the ergodic
strategy also exhibits greater robustness to noise, maintaining a
consistent performance advantage across all evaluated conditions.
Overall, these findings confirm the benefit of ergodic exploration
in balancing detection reliability, reconstruction accuracy, and
execution efficiency once suspicious regions have been identified.

\vspace{0.6em}

% ---------- D. Segmentation ----------
\noindent\textit{D.~Segmentation}

\vspace{0.2em}

\noindent
The spatial accuracy of the proposed method was evaluated based on
its ability to segment stiff regions within the reconstructed
stiffness maps and compared against the BO. To this end, we designed
two simulated stiffness distributions: the first one with two
inclusions, one irregularly shaped (Cluster 1) and one circular
(Cluster 2), while the second distribution consisted of a single
donut-shaped inclusion (Cluster 3), as

% ---------- Bottom license strip ----------
\vfill
\hrule
\vspace{0.3em}
\begin{center}
{\footnotesize\itshape
This work is licensed under a Creative Commons Attribution 4.0
License. For more information, see
\url{https://creativecommons.org/licenses/by/4.0/}\par}
\end{center}

\end{document}
